% Figura: problema 4.7. Electrete cilíndrico coaxial con P = (k r/2) u_r.
\begin{figure}[t]
\centering
\begin{subfigure}[b]{0.27\textwidth}\centering
\begin{tikzpicture}[scale=0.8]
  % cilindro hueco en perspectiva
  \fill[dielec] (-1.3,0) -- (-1.3,-3.0) arc (180:360:1.3 and 0.38) -- (1.3,0) arc (0:180:1.3 and 0.38);
  \fill[dielec] (-1.3,0) arc (180:360:1.3 and 0.38) -- (1.3,0) arc (0:180:1.3 and 0.38);
  \fill[dielecB] (0,0) ellipse (1.3 and 0.38);
  \fill[fondo] (0,0) ellipse (0.55 and 0.16);
  \draw[borde] (0,0) ellipse (1.3 and 0.38) (0,0) ellipse (0.55 and 0.16);
  \draw[borde] (-1.3,0) -- (-1.3,-3.0) arc (180:360:1.3 and 0.38) -- (1.3,0);
  \draw[texto!50!fondo, dashed, line width=0.4pt] (-0.55,0) -- (-0.55,-3.0) (0.55,0) -- (0.55,-3.0);
  \draw[texto!50!fondo, dashed, line width=0.4pt] (-0.55,-3.0) arc (180:360:0.55 and 0.16);
  \foreach \t in {20,70,...,340}
    \draw[campoP, line width=0.7pt, -{Stealth[length=1.3mm]}] ({0.6*cos(\t)},{0.175*sin(\t)}) -- ({1.22*cos(\t)},{0.355*sin(\t)});
  \draw[cota] (1.65,0) -- (1.65,-3.0) node[midway, right, etiq] {$\ell$};
  \node[etiq] at (-0.25,0.52) {$a$};
  \node[etiq] at (1.0,0.62) {$b$};
\end{tikzpicture}
\caption{Electrete coaxial}
\end{subfigure}\hfill
\begin{subfigure}[b]{0.27\textwidth}\centering
\begin{tikzpicture}[scale=0.8]
  \fill[dielec] (0,0) circle (1.7);
  \fill[fondo] (0,0) circle (0.8);
  \draw[borde] (0,0) circle (1.7) (0,0) circle (0.8);
  \foreach \t in {0,45,...,315}
    \draw[campoP, line width=0.8pt, -{Stealth[length=1.5mm]}] (\t:0.9) -- (\t:{0.9+0.62*1.6/1.7});
  \draw[vecn] (135:0.8) -- (135:0.3) node[right=-1pt, etiq] {$\nn$};
  \draw[vecn] (1.7,0) -- (2.3,0) node[above, etiq] {$\nn$};
  \node[etiq, text=campoP] at (62:2.0) {$\vP$};
  \node[etiq, font=\footnotesize] at (0,-0.38) {$\vE=\vcero$};
\end{tikzpicture}
\caption{Corte transversal}
\end{subfigure}\hfill
\begin{subfigure}[b]{0.42\textwidth}\centering
\begin{tikzpicture}
\begin{axis}[width=5.8cm, height=4.3cm, axis lines=middle, xmin=0, xmax=3.1, ymin=-1.15, ymax=1.15,
  xtick={1,2}, xticklabels={$a$,$b$}, ytick=\empty, xlabel={$r$},
  xlabel style={at={(ticklabel* cs:1)}, anchor=north east}, clip=false,
  legend style={at={(0.5,-0.08)}, anchor=north, legend columns=2, font=\tiny, /tikz/every even column/.append style={column sep=4pt}}]
  \addplot[curvaP, domain=1:2, samples=2] {x/2}; \addlegendentry{$P=kr/2$}
  \addplot[curvaE, domain=1:2, samples=2] {-x/2}; \addlegendentry{$\eo E_r=-kr/2$}
  \addplot[curvaD] coordinates {(0,0.02) (3.1,0.02)}; \addlegendentry{$D_r=0$}
  \addplot[curvaP, forget plot] coordinates {(1,0) (1,0.5)};
  \addplot[curvaP, forget plot] coordinates {(2,0) (2,1)};
  \draw[equipot] (axis cs:1,-1.15) -- (axis cs:1,1.15) (axis cs:2,-1.15) -- (axis cs:2,1.15);
  \node[etiqt, anchor=west] at (axis cs:2.05,1.0) {$\sigmaP(b)=\tfrac{kb}{2}$};
  \node[etiqt, anchor=west] at (axis cs:1.03,-0.22) {$\sigmaP(a)=-\tfrac{ka}{2}$};
  \node[etiqt] at (axis cs:1.32,-1.05) {$\rhoP=-k$};
\end{axis}
\end{tikzpicture}
\caption{Perfiles radiales}
\end{subfigure}
\caption{\Cref{prob:electrete-coaxial}. (a), (b) Electrete limitado por dos cilindros coaxiales de radios $a<b$, con polarización
radial creciente. (c) Hay carga ligada de volumen uniforme, $\rhoP=-k$, y de superficie en ambas caras;
su total por unidad de longitud es nulo. Como no hay carga libre y $\vP$ es irrotacional, $\vD=\vcero$ y
$\vE=-\vP/\eo$ dentro del material; fuera de él, el campo es nulo.}
\label{fig:p47}
\end{figure}
